\documentclass[11pt]{book}
\usepackage[paperwidth=145mm,paperheight=200mm,top=18mm,bottom=20mm,inner=20mm,outer=16mm]{geometry}
\usepackage{brumfiel}
% figures03.tex -- Chapter 3 ("Concerning Lines") figures redrawn in TikZ.
%
% Fig 3-1 / 3-2 geometry is DIGITIZED from scan04.pdf p4 (book p.44): the six
% straight lines of Fig 3-1 were recovered by a Hough fit over the thresholded
% scan crop, and Fig 3-2 is the same six lines carrying the printed hand-wobble
% (measured perpendicular deviation 8-15 px at 170 dpi ~= 0.04-0.075 cm).
% Remaining figures are reconstructions at printed proportion; constrained
% points are placed by construction (never by eye) -- see tools/constraints/ch03.py.

% ============================================================ Fig 3-1, 3-2
% Six lines. L1 baseline carries A,B; L3 carries A',B'.
% AB and A'B' are drawn the SAME length in 3-1 (the text's whole point).
\newcommand{\FIGIIIONETWO}{\begin{bkfigure}[\linewidth]
\begin{minipage}{0.48\linewidth}\centering
\begin{tikzpicture}[scale=0.62]
  % digitized endpoints (scan px / 200, y flipped)
  \coordinate (p1) at (0.18,0.985);  \coordinate (q1) at (8.69,0.80);
  \coordinate (p2) at (0.905,3.195); \coordinate (q2) at (7.945,2.315);
  \coordinate (p3) at (4.06,4.075);  \coordinate (q3) at (7.60,0.81);
  \coordinate (p4) at (4.64,4.08);   \coordinate (q4) at (1.81,0.625);
  \coordinate (p5) at (2.005,3.565); \coordinate (q5) at (3.25,0.005);
  \coordinate (p6) at (7.35,3.64);   \coordinate (q6) at (5.74,0.765);
  \foreach \a/\b in {p1/q1,p2/q2,p3/q3,p4/q4,p5/q5,p6/q6} \draw[fig] (\a)--(\b);
  % A,B on the baseline; A',B' on L3 -- equal separations by construction
  \coordinate (A) at ($(p1)!0.085!(q1)$);  \coordinate (B) at ($(p1)!0.150!(q1)$);
  % A'B' drawn the SAME LENGTH as AB: t-step scaled by |L1|/|L3|.
  % Both slid down L3 past its crossing with L6 (t=0.658) so the labels have
  % clear air on the right-hand side of the line.
  \coordinate (Ap) at ($(p3)!0.83489!(q3)$); \coordinate (Bp) at ($(p3)!0.72000!(q3)$);
  \foreach \p in {A,B,Ap,Bp} {\dt{\p}}
  \node[vlab,above] at (A) {$A$};   \node[vlab,above] at (B) {$B$};
  \node[vlab,right,yshift=4pt] at (Ap) {$A'$};
  \node[vlab] at ($(Bp)+(47.3:0.45)$) {$B'$};
\end{tikzpicture}
\figcap{3--1}
\end{minipage}\hfill
\begin{minipage}{0.48\linewidth}\centering
\begin{tikzpicture}[scale=0.62]
  % same six lines seen "through uneven glass": wobble traced from the scan
  \draw[fig] plot[smooth,tension=0.85] coordinates
    {(0.245,0.745)(1.30,0.79)(2.35,0.755)(3.40,0.815)(4.45,0.80)(5.50,0.86)(6.45,0.90)(7.40,0.945)};
  \draw[fig] plot[smooth,tension=0.85] coordinates
    {(1.005,2.98)(1.90,2.90)(2.80,2.90)(3.70,2.79)(4.60,2.68)(5.50,2.62)(6.45,2.46)(7.39,2.355)};
  \draw[fig] plot[smooth,tension=0.85] coordinates
    {(4.05,3.87)(4.52,3.42)(4.95,2.98)(5.40,2.55)(5.85,2.09)(6.30,1.68)(6.75,1.30)(7.185,0.905)};
  \draw[fig] plot[smooth,tension=0.85] coordinates
    {(4.515,3.98)(4.14,3.53)(3.73,3.04)(3.38,2.55)(2.99,2.03)(2.63,1.53)(2.29,1.03)(1.945,0.515)};
  \draw[fig] plot[smooth,tension=0.85] coordinates
    {(1.97,3.355)(2.11,2.92)(2.28,2.49)(2.40,2.06)(2.56,1.64)(2.72,1.21)(2.93,0.79)(3.09,0.37)};
  \draw[fig] plot[smooth,tension=0.85] coordinates
    {(6.835,3.555)(6.62,3.11)(6.42,2.66)(6.24,2.20)(6.02,1.75)(5.83,1.30)(5.63,0.85)(5.43,0.40)};
  % same two pairs as Fig 3-1, at the same fractions along their curves
  \coordinate (A) at (0.83,0.775);  \coordinate (B) at (1.28,0.79);
  \coordinate (Ap) at (6.681,1.359); \coordinate (Bp) at (6.319,1.664);
  \foreach \p in {A,B,Ap,Bp} {\dt{\p}}
  \node[vlab,above] at (A) {$A$};   \node[vlab,above] at (B) {$B$};
  \node[vlab,right,yshift=4pt] at (Ap) {$A'$};
  \node[vlab] at ($(Bp)+(49.8:0.50)$) {$B'$};
\end{tikzpicture}
\figcap{3--2}
\end{minipage}
\end{bkfigure}}

% ======================================================= Fig 3-3, 3-4, 3-5
\newcommand{\FIGIIITHREEFOURFIVE}{\begin{bkfigure}[\linewidth]
\begin{minipage}{0.30\linewidth}\centering
\begin{tikzpicture}[scale=0.62]
  \draw[fig] (0,0) -- (5.0,0.55); \node[vlab,right] at (5.0,0.55) {$l$};
  \coordinate (P) at (2.0,0.22); \dt{P}
  \node[vlab,above] at (P) {$P$};
\end{tikzpicture}
\figcap{3--3}
\end{minipage}\hfill
\begin{minipage}{0.30\linewidth}\centering
\begin{tikzpicture}[scale=0.62]
  \coordinate (la) at (0,0.10); \coordinate (lb) at (4.8,1.42);
  \coordinate (ma) at (0,1.55); \coordinate (mb) at (4.8,0.02);
  \draw[fig] (la) -- (lb); \node[vlab,right] at (lb) {$l$};
  \draw[fig] (ma) -- (mb); \node[vlab,right] at (mb) {$m$};
  \coordinate (P) at (intersection of la--lb and ma--mb);
  \dt{P} \node[vlab,below] at (P) {$P$};
\end{tikzpicture}
\figcap{3--4}
\end{minipage}\hfill
\begin{minipage}{0.30\linewidth}\centering
\begin{tikzpicture}[scale=0.62]
  \coordinate (A) at (2.35,2.55); \coordinate (B) at (0.75,0.32);
  \coordinate (C) at (4.25,0.32);
  \draw[fig] ($(A)!-0.16!(B)$) -- ($(B)!-0.16!(A)$);
  \draw[fig] ($(A)!-0.16!(C)$) -- ($(C)!-0.16!(A)$);
  \draw[fig] ($(B)!-0.20!(C)$) -- ($(C)!-0.20!(B)$);
  % labels on the bisector of the free wedge at each vertex (the book puts
  % A above the crossing, B and C just inside the base)
  \node[vlab] at ($(A)+(92.5:0.75)$) {$A$};
  \node[vlab] at ($(B)+(27:0.90)$)   {$B$};
  \node[vlab] at ($(C)+(155:0.95)$)  {$C$};
\end{tikzpicture}
\figcap{3--5}
\end{minipage}
\end{bkfigure}}

% ============================================================ Fig 3-6, 3-7
% 3-6: two CURVED "lines" l and m meeting at A and B (the impossible picture).
\newcommand{\FIGIIISIXSEVEN}{\begin{bkfigure}[\linewidth]
\begin{minipage}{0.48\linewidth}\centering
\begin{tikzpicture}[scale=0.714]
  \coordinate (A) at (0.45,0.55); \coordinate (B) at (4.65,1.15);
  \draw[fig] (A) .. controls (1.7,1.55) and (3.3,1.75) .. (B);
  \draw[fig] (A) .. controls (1.7,0.10) and (3.3,0.35) .. (B);
  % BOTH curves run on past BOTH meeting points, so A and B each read as a
  % crossing (as in the book), not as a cusp
  \draw[fig] ($(A)!-0.09!(1.7,1.55)$) -- (A);
  \draw[fig] ($(A)!-0.09!(1.7,0.10)$) -- (A);
  \draw[fig] (B) -- ($(B)!-0.09!(3.3,1.75)$);
  \draw[fig] (B) -- ($(B)!-0.09!(3.3,0.35)$);
  \dt{A} \dt{B}
  \node[vlab,left] at (A) {$A$}; \node[vlab] at ($(B)+(100:0.50)$) {$B$};
  \node[vlab,above right] at (3.92,1.50) {$l$};
  \node[vlab,below right] at (4.55,0.88) {$m$};
\end{tikzpicture}
\figcap{3--6}
\end{minipage}\hfill
\begin{minipage}{0.48\linewidth}\centering
\begin{tikzpicture}[scale=0.714]
  \coordinate (C) at (2.55,2.60); \coordinate (A) at (0.55,0.35);
  \coordinate (B) at (4.35,0.62);
  \draw[fig] ($(A)!-0.18!(B)$) -- ($(B)!-0.16!(A)$);
  \draw[fig] ($(C)!-0.20!(A)$) -- ($(A)!-0.10!(C)$);
  \draw[fig] ($(C)!-0.20!(B)$) -- ($(B)!-0.10!(C)$);
  % m = line BC (extends up-LEFT above C); n = line CA (extends up-RIGHT).
  % C sits in the wedge ABOVE the crossing, between m and n, as in the book.
  \node[vlab] at ($(C)+(90:0.72)$) {$C$};
  \node[vlab,above left] at (A) {$A$}; \node[vlab,below right] at (B) {$B$};
  \node[vlab,left]  at ($(C)!-0.20!(B)$) {$m$};
  \node[vlab,right] at ($(C)!-0.20!(A)$) {$n$};
  \node[vlab,right] at ($(B)!-0.16!(A)$) {$l$};
\end{tikzpicture}
\figcap{3--7}
\end{minipage}
\end{bkfigure}}

% ============================================================ Fig 3-8, 3-9
\newcommand{\FIGIIIEIGHT}{\begin{bkfigure}[0.62\linewidth]
\begin{tikzpicture}[scale=0.806]
  \coordinate (P) at (0.40,0.35); \coordinate (Q) at (4.55,1.35);
  \draw[fig] (P) .. controls (1.7,1.35) and (3.2,1.70) .. (Q);
  \draw[fig] (P) .. controls (1.8,0.05) and (3.4,0.55) .. (Q);
  \draw[fig] ($(P)!-0.10!(1.7,1.35)$) -- (P);
  \draw[fig] ($(P)!-0.10!(1.8,0.05)$) -- (P);
  \draw[fig] (Q) -- ($(Q)!-0.10!(3.2,1.70)$);
  \draw[fig] (Q) -- ($(Q)!-0.10!(3.4,0.55)$);
  \dt{P} \dt{Q}
  \node[vlab,left] at (P) {$P$}; \node[vlab,above left] at (Q) {$Q$};
  \node[vlab,above right] at (4.30,1.60) {$a$};
  \node[vlab,below right] at (4.55,1.12) {$b$};
\end{tikzpicture}
\figcap{3--8}
\end{bkfigure}}

% 3-9: figure-eight -- two curved "lines" l and m crossing twice, A..E marked.
\newcommand{\FIGIIININE}{\begin{bkfigure}[0.62\linewidth]
\begin{tikzpicture}[scale=0.806]
  \coordinate (A) at (0.30,1.05); \coordinate (B) at (2.45,1.05);
  \coordinate (C) at (4.60,0.92);
  % D and E sit ON the right lobe: Bezier midpoints of the two arcs B..C
  \coordinate (D) at (3.5625,1.74625);
  \coordinate (E) at (3.5625,0.32125);
  \draw[fig] (A) .. controls (0.95,2.05) and (1.85,2.00) .. (B)
                 .. controls (3.05,2.05) and (4.10,1.95) .. (C);
  \draw[fig] (A) .. controls (0.95,0.05) and (1.85,0.10) .. (B)
                 .. controls (3.05,0.05) and (4.10,0.15) .. (C);
  \draw[fig] ($(A)!-0.12!(0.95,2.05)$) -- (A);
  \draw[fig] ($(A)!-0.12!(0.95,0.05)$) -- (A);
  \draw[fig] (C) -- ($(C)!-0.14!(4.10,1.95)$);
  \draw[fig] (C) -- ($(C)!-0.14!(4.10,0.15)$);
  \foreach \p in {A,B,C,D,E} {\dt{\p}}
  \node[vlab,above left] at (A) {$A$}; \node[vlab] at ($(B)+(90:0.58)$) {$B$};
  \node[vlab,below] at (C) {$C$};
  \node[vlab,above] at (D) {$D$}; \node[vlab,below] at (E) {$E$};
  \node[vlab,above right] at (4.80,1.36) {$l$};
  \node[vlab,right] at (5.02,0.62) {$m$};
\end{tikzpicture}
\figcap{3--9}
\end{bkfigure}}

% =============================================================== Fig 3-10
\newcommand{\FIGIIITEN}{\begin{bkfigure}[0.55\linewidth]
\begin{tikzpicture}[scale=0.806]
  \coordinate (A) at (2.45,2.85); \coordinate (B) at (0.35,0.30);
  \coordinate (C) at (4.55,0.55);
  \draw[fig] (A)--(B)--(C)--cycle;
  % D on AB, E on AC; the drawn line DE extended both ways
  \coordinate (D) at ($(A)!0.52!(B)$); \coordinate (E) at ($(A)!0.52!(C)$);
  \draw[fig] ($(D)!-0.30!(E)$) -- ($(E)!-0.22!(D)$);
  \foreach \p in {A,B,C,D,E} {\dt{\p}}
  \node[vlab,above] at (A) {$A$}; \node[vlab,below left] at (B) {$B$};
  \node[vlab,right] at (C) {$C$};
  \node[vlab,above left] at (D) {$D$}; \node[vlab,above right] at (E) {$E$};
\end{tikzpicture}
\figcap{3--10}
\end{bkfigure}}

% ========================================================= Fig 3-11, 3-12
\newcommand{\FIGIIIELEVENTWELVE}{\begin{bkfigure}[\linewidth]
\begin{minipage}{0.48\linewidth}\centering
\begin{tikzpicture}[scale=0.762]
  \draw[fig] (0,0) -- (4.6,0.62);
  \coordinate (A) at (0.85,0.115); \coordinate (B) at (2.95,0.398);
  \coordinate (C) at (3.55,0.478);
  \foreach \p in {A,B,C} {\dt{\p}}
  \node[vlab,above] at (A) {$A$}; \node[vlab,above] at (B) {$B$};
  \node[vlab,above] at (C) {$C$};
\end{tikzpicture}
\figcap{3--11}
\end{minipage}\hfill
\begin{minipage}{0.48\linewidth}\centering
\begin{tikzpicture}[scale=0.762]
  \draw[fig] (0,0) -- (4.6,0.62);
  \coordinate (A) at (0.60,0.081); \coordinate (D) at (2.05,0.276);
  \coordinate (B) at (2.60,0.350); \coordinate (C) at (4.05,0.546);
  \foreach \p in {A,D,B,C} {\dt{\p}}
  \node[vlab,above] at (A) {$A$}; \node[vlab,above] at (D) {$D$};
  \node[vlab,above] at (B) {$B$}; \node[vlab,above] at (C) {$C$};
\end{tikzpicture}
\figcap{3--12}
\end{minipage}
\end{bkfigure}}

% ========================================================= Fig 3-13, 3-14
\newcommand{\FIGIIITHIRTEENFOURTEEN}{\begin{bkfigure}[\linewidth]
\begin{minipage}{0.48\linewidth}\centering
\begin{tikzpicture}[scale=0.741]
  \draw[fig] (0,0) -- (4.6,1.55); \node[vlab,right] at (4.6,1.55) {$l$};
  \coordinate (P) at (0.55,1.95); \coordinate (Q) at (2.25,1.95);
  \draw[fig] (P) -- (Q);
  \dt{P} \dt{Q}
  \node[vlab,above] at (P) {$P$}; \node[vlab,above] at (Q) {$Q$};
  \node[slab,below] at ($(P)!0.5!(Q)$) {Same side};
\end{tikzpicture}
\figcap{3--13}
\end{minipage}\hfill
\begin{minipage}{0.48\linewidth}\centering
\begin{tikzpicture}[scale=0.741]
  \coordinate (P) at (1.25,2.10); \coordinate (Q) at (2.85,0.20);
  \coordinate (la) at (0.30,0.05); \coordinate (lb) at (4.30,2.30);
  \draw[fig] (la) -- (lb); \node[vlab,right] at (lb) {$l$};
  \draw[fig] (P) -- (Q);
  % R is where PQ meets l -- by construction, never by eye
  \coordinate (R) at (intersection of la--lb and P--Q);
  \foreach \p in {P,Q,R} {\dt{\p}}
  \node[vlab,above] at (P) {$P$}; \node[vlab,below] at (Q) {$Q$};
  \node[vlab] at ($(R)+(170:0.80)$) {$R$};
  \node[slab,right,align=left] at (3.05,1.14) {Opposite\\sides};
\end{tikzpicture}
\figcap{3--14}
\end{minipage}
\end{bkfigure}}

% =============================================================== Fig 3-15
% Halftone-style line illustration of three men standing in a line: PLATE.
\newcommand{\FIGIIIFIFTEEN}{\begin{bkfigure}[0.72\linewidth]
\setlength{\fboxsep}{10pt}%
\fbox{\begin{minipage}{0.86\linewidth}\centering\footnotesize
\textsc{Plate.}\ Illustration: three men standing in a line, marked
$A$, $B$, $C$, with the sight line from $A$ blocked at $B$.\\[2pt]
\emph{Original artwork -- not redrawn.}
\end{minipage}}
\figcap{3--15}
\end{bkfigure}}

% ========================================================= Fig 3-16, 3-17
\newcommand{\FIGIIISIXTEENSEVENTEEN}{\begin{bkfigure}[\linewidth]
\begin{minipage}{0.48\linewidth}\centering
\begin{tikzpicture}[scale=0.768]
  \coordinate (A) at (0.20,0.15); \coordinate (B) at (3.55,2.20);
  \draw[fig] (A) -- (B);
  \coordinate (P) at ($(A)!0.46!(B)$);
  \foreach \p in {A,B,P} {\dt{\p}}
  \node[vlab,below left] at (A) {$A$}; \node[vlab,above] at (B) {$B$};
  \node[vlab,above left] at (P) {$P$};
\end{tikzpicture}
\figcap{3--16}
\end{minipage}\hfill
\begin{minipage}{0.48\linewidth}\centering
\begin{tikzpicture}[scale=0.768]
  \coordinate (A) at (2.05,2.45); \coordinate (B) at (0.25,0.20);
  \coordinate (C) at (3.75,0.05);
  \draw[fig] (A)--(B)--(C)--cycle;
  \node[vlab,above] at (A) {$A$}; \node[vlab,left] at (B) {$B$};
  \node[vlab,below right] at (C) {$C$};
\end{tikzpicture}
\figcap{3--17}
\end{minipage}
\end{bkfigure}}

% =============================================================== Fig 3-18
% l meets side AB and misses BC; C is the apex. l passes through no vertex.
\newcommand{\FIGIIIEIGHTEEN}{\begin{bkfigure}[0.55\linewidth]
\begin{tikzpicture}[scale=0.806]
  \coordinate (C) at (1.85,2.60); \coordinate (B) at (0.20,0.30);
  \coordinate (A) at (4.30,0.35);
  \draw[fig] (C)--(B)--(A)--cycle;
  % l: enters through AB (the base) and leaves through AC
  \coordinate (X) at ($(B)!0.42!(A)$);      % on AB (base)
  \coordinate (Y) at ($(C)!0.52!(A)$);      % on AC
  % l runs from just below the base up past AC to about C's height, as printed
  \coordinate (ltop) at ($(Y)!-1.05!(X)$);
  \draw[fig] ($(X)!-0.35!(Y)$) -- (ltop);
  \node[vlab,above] at (C) {$C$}; \node[vlab,left] at (B) {$B$};
  \node[vlab,right] at (A) {$A$};
  \node[vlab,above right] at (ltop) {$l$};
\end{tikzpicture}
\figcap{3--18}
\end{bkfigure}}

% =============================================================== Fig 3-19
\newcommand{\FIGIIININETEEN}{\begin{bkfigure}[0.62\linewidth]
\begin{tikzpicture}[scale=0.806]
  \draw[fig] (0,0) -- (5.4,0); \node[vlab,right] at (5.4,0) {$l$};
  \coordinate (O) at (2.55,0); \dt{O}
  \node[vlab,below] at (O) {$O$};
  \draw[fig,line width=0.5pt] (0.35,0.30) -- (0.35,0.46) -- (2.45,0.46) -- (2.45,0.30);
  \draw[fig,line width=0.5pt] (2.65,0.30) -- (2.65,0.46) -- (4.95,0.46) -- (4.95,0.30);
  \node[slab,above] at (1.40,0.46) {One side};
  \node[slab,above] at (3.80,0.46) {Other side};
\end{tikzpicture}
\figcap{3--19}
\end{bkfigure}}

% ========================================================= Fig 3-20, 3-21
\newcommand{\FIGIIITWENTYTWENTYONE}{\begin{bkfigure}[\linewidth]
\begin{minipage}{0.48\linewidth}\centering
\begin{tikzpicture}[scale=0.751]
  \coordinate (la) at (0,0.55);     \coordinate (lb) at (4.35,0.55);
  \coordinate (ma) at (1.05,-0.45); \coordinate (mb) at (2.95,2.40);
  \draw[fig] (la) -- (lb); \node[vlab,right] at (lb) {$l$};
  \draw[fig] (ma) -- (mb); \node[vlab,above] at (mb) {$m$};
  % O is the point l and m have in common; A is a point of m, not of l
  \coordinate (O) at (intersection of la--lb and ma--mb);
  \coordinate (A) at ($(ma)!0.79!(mb)$);
  \dt{O} \dt{A}
  \node[vlab] at ($(O)+(298:0.55)$) {$O$}; \node[vlab,right] at (A) {$A$};
\end{tikzpicture}
\figcap{3--20}
\end{minipage}\hfill
\begin{minipage}{0.48\linewidth}\centering
\begin{tikzpicture}[scale=0.751]
  \coordinate (la) at (0,0.55);     \coordinate (lb) at (4.35,0.55);
  \coordinate (ma) at (1.05,-0.45); \coordinate (mb) at (2.95,2.40);
  \draw[fig] (la) -- (lb); \node[vlab,right] at (lb) {$l$};
  \draw[fig] (ma) -- (mb); \node[vlab,above] at (mb) {$m$};
  \coordinate (O) at (intersection of la--lb and ma--mb);
  \coordinate (A) at ($(ma)!0.79!(mb)$);
  % P and Q are points of l on opposite sides of O:  POQ
  \coordinate (Q) at ($(la)!0.138!(lb)$); \coordinate (P) at ($(la)!0.655!(lb)$);
  \foreach \p in {O,A,P,Q} {\dt{\p}}
  \node[vlab] at ($(O)+(298:0.55)$) {$O$}; \node[vlab,right] at (A) {$A$};
  \node[vlab,below] at (Q) {$Q$}; \node[vlab,below] at (P) {$P$};
\end{tikzpicture}
\figcap{3--21}
\end{minipage}
\end{bkfigure}}

% ========================================================= Fig 3-22, 3-23
\newcommand{\FIGIIITWENTYTWOTWENTYTHREE}{\begin{bkfigure}[\linewidth]
\begin{minipage}{0.48\linewidth}\centering
\begin{tikzpicture}[scale=0.755]
  \coordinate (O) at (1.05,0.35); \coordinate (lend) at (4.35,1.32);
  % the part of l that is NOT in the ray, dashed: the exact opposite ray
  \draw[fig,dashed,dash pattern=on 3pt off 2.4pt] ($(O)!-0.318!(lend)$) -- (O);
  \draw[fig] (O) -- (lend); \node[vlab,right] at (lend) {$l$};
  \coordinate (P) at ($(O)!0.394!(lend)$);
  \dt{O} \dt{P}
  \node[vlab,above] at (O) {$O$}; \node[vlab,above] at (P) {$P$};
  \node[slab,below] at (1.90,0.52) {Ray};
\end{tikzpicture}
\figcap{3--22}
\end{minipage}\hfill
\begin{minipage}{0.48\linewidth}\centering
\begin{tikzpicture}[scale=0.755]
  \coordinate (P) at (2.05,0.30);
  \draw[fig] (0.15,0.30) -- (4.35,0.30); \node[vlab,right] at (4.35,0.30) {$l$};
  \draw[fig] (P) -- (3.35,2.55); \node[vlab,above] at (3.35,2.55) {$r$};
  \dt{P} \node[vlab,above left] at (P) {$P$};
\end{tikzpicture}
\figcap{3--23}
\end{minipage}
\end{bkfigure}}

% ========================================================= Fig 3-24, 3-25
\newcommand{\FIGIIITWENTYFOURTWENTYFIVE}{\begin{bkfigure}[\linewidth]
\begin{minipage}{0.52\linewidth}\centering
\begin{tikzpicture}[scale=0.780]
  % (i) angle at O, upper left
  \coordinate (O1) at (0.30,2.05);
  \draw[fig] (O1) -- (1.85,2.95); \draw[fig] (O1) -- (1.85,2.05);
  \dt{O1} \node[vlab,left] at (O1) {$O$};
  % (ii) angle at O, lower left
  \coordinate (O2) at (0.62,1.35);
  \draw[fig] (O2) -- (2.05,1.05); \draw[fig] (O2) -- (-0.15,0.20);
  \dt{O2} \node[vlab,left] at (O2) {$O$};
  % (iii) angle AOB, right
  \coordinate (O3) at (3.55,2.20); \coordinate (A3) at (2.55,3.10);
  \coordinate (B3) at (2.85,1.25);
  \draw[fig] (O3) -- (A3); \draw[fig] (O3) -- (B3);
  \dt{O3} \dt{A3} \dt{B3}
  \node[vlab,right] at (O3) {$O$}; \node[vlab,above] at (A3) {$A$};
  \node[vlab,below] at (B3) {$B$};
  % (iv) straight angle
  \draw[fig] (0.85,0.00) -- (3.35,0.00);
  \coordinate (O4) at (2.05,0.00); \dt{O4}
  \node[vlab,below] at (O4) {$O$};
\end{tikzpicture}
\figcap{3--24}
\end{minipage}\hfill
\begin{minipage}{0.44\linewidth}\centering
\begin{tikzpicture}[scale=0.780]
  % (a) the arrowheads are on the ENDS OF THE ARC, pointing at the two sides;
  % the sides themselves carry no arrows.  Verified on scans/scan04.pdf p18
  % (book p.54), where both arrowheads are plainly on the arc.
  \coordinate (Oa) at (0.15,2.55);
  \draw[fig] (Oa) -- (1.90,3.15);
  \draw[fig] (Oa) -- (1.75,2.10);
  % every arc is centred on its vertex and stops exactly on the two sides
  \draw[fig,line width=0.5pt,<->,>=stealth]
    ($(Oa)+(18.93:1.05)$) arc (18.93:-15.71:1.05);
  \node[slab] at ($(Oa)+(1.61:1.30)$) {1};
  \node[slab,right] at (2.25,2.55) {(a)};
  % (b) plain sides
  \coordinate (Ob) at (0.15,1.35);
  \draw[fig] (Ob) -- (1.90,1.95); \draw[fig] (Ob) -- (1.75,0.90);
  \draw[fig,line width=0.5pt] ($(Ob)+(18.93:1.05)$) arc (18.93:-15.71:1.05);
  \node[slab] at ($(Ob)+(1.61:1.30)$) {1};
  \node[slab,right] at (2.25,1.35) {(b)};
  % (c) obtuse
  \coordinate (Oc) at (0.35,0.10);
  \draw[fig] (Oc) -- (2.10,0.10); \draw[fig] (Oc) -- (-0.25,0.95);
  \draw[fig,line width=0.5pt] ($(Oc)+(0:0.90)$) arc (0:125.18:0.90);
  \node[slab] at ($(Oc)+(62.59:0.58)$) {2};
  \node[slab,right] at (2.25,0.10) {(c)};
\end{tikzpicture}
\figcap{3--25}
\end{minipage}
\end{bkfigure}}

% =============================================================== Fig 3-26
% Two lines l, m crossing at O; four rays l1,l2,m1,m2; angles 1,2,3,4.
\newcommand{\FIGIIITWENTYSIX}{\begin{bkfigure}[0.66\linewidth]
\begin{tikzpicture}[scale=0.910]
  \coordinate (L1) at (4.85,2.05); \coordinate (L2) at (0.35,0.15);
  \coordinate (M1) at (4.85,0.35); \coordinate (M2) at (0.35,1.95);
  \draw[fig] (L2) -- (L1); \draw[fig] (M2) -- (M1);
  \coordinate (O) at (intersection of L2--L1 and M2--M1);
  \dt{O}
  \node[vlab,right] at (L1) {$l_1$}; \node[vlab,left]  at (L2) {$l$};
  \node[vlab,right] at (M1) {$m_1$}; \node[vlab,left]  at (M2) {$m$};
  % ray labels offset perpendicular to their own line, so they sit beside it
  \node[vlab] at ($($(M2)!0.28!(M1)$)+(70.42:0.42)$) {$m_2$};
  \node[vlab] at ($($(L2)!0.30!(L1)$)+(-67.10:0.42)$) {$l_2$};
  % the four arcs are centred on O and stop exactly on the four rays
  \draw[fig,line width=0.5pt] ($(O)+(160.44:0.62)$) arc (160.44:202.89:0.62);
  \draw[fig,line width=0.5pt] ($(O)+(-19.58:0.62)$) arc (-19.58:22.90:0.62);
  \draw[fig,line width=0.5pt] ($(O)+(22.90:0.62)$)  arc (22.90:160.44:0.62);
  \draw[fig,line width=0.5pt] ($(O)+(202.89:0.62)$) arc (202.89:340.42:0.62);
  % angle numerals on each region's bisector, clear of the arcs
  \node[slab] at ($(O)+(181.67:0.95)$) {1};
  \node[slab] at ($(O)+(1.66:0.95)$)   {2};
  \node[slab] at ($(O)+(91.67:0.95)$)  {3};
  \node[slab] at ($(O)+(271.66:0.95)$) {4};
  % the book sets O just below-left of the crossing, inside the ring of arcs;
  % at this type size the label padding will not clear the two lines there, so
  % it stays outside the ring, below-right (see ch03-UNCERTAIN.md)
  \node[vlab] at ($(O)+(303:1.00)$) {$O$};
\end{tikzpicture}
\figcap{3--26}
\end{bkfigure}}

% =============================================================== Fig 3-27
% Three drawings of a pair of angles sharing a vertex and one side.
\newcommand{\FIGIIITWENTYSEVEN}{\begin{bkfigure}[0.82\linewidth]
\begin{tikzpicture}[scale=0.780]
  % (i) narrow pair
  % the narrowest of the three pairs; opened from 18/24 deg to 23/25 so the
  % numerals clear both rays at this type size (the book's own fan is tight)
  \coordinate (O) at (0.15,0.55);
  \draw[fig] (O) -- (1.891,1.724); \draw[fig] (O) -- (2.261,0.960);
  \draw[fig] (O) -- (2.091,0.066);
  \draw[fig,line width=0.5pt] ($(O)+(34:1.25)$) arc (34:11:1.25);
  \draw[fig,line width=0.5pt] ($(O)+(11:1.25)$) arc (11:-14:1.25);
  \node[slab] at ($(O)+(22.5:1.62)$) {1};
  \node[slab] at ($(O)+(-1.5:1.62)$) {2};
\end{tikzpicture}\hspace{1.6em}
\begin{tikzpicture}[scale=0.780]
  % (ii) the book's middle drawing: ONE ray straight up, TWO rays down-left and
  % down-right (an inverted Y), the vertical being the common side.  The two
  % arcs bulge out either side of it -- measured off the printed figure.
  \coordinate (O) at (1.15,1.35);
  \draw[fig] (O) -- (1.15,2.75);            % common side, straight up
  \draw[fig] (O) -- (0.30,0.25);            % noncommon side, down-left
  \draw[fig] (O) -- (2.05,0.25);            % noncommon side, down-right
  \draw[fig,line width=0.5pt] ($(O)+(232.30:0.80)$) arc (232.30:90:0.80);
  \draw[fig,line width=0.5pt] ($(O)+(90:0.80)$) arc (90:-50.71:0.80);
  \node[slab] at ($(O)+(161.15:1.14)$) {1};
  \node[slab] at ($(O)+(19.65:1.14)$)  {2};
\end{tikzpicture}\hspace{1.6em}
\begin{tikzpicture}[scale=0.780]
  % (iii) wide pair over a base line
  \coordinate (O) at (1.25,0.25);
  \draw[fig] (O) -- (0.05,1.35); \draw[fig] (O) -- (1.55,2.05);
  \draw[fig] (O) -- (2.85,0.90);
  \draw[fig,line width=0.5pt] ($(O)+(137.48:1.35)$) arc (137.48:80.54:1.35);
  \draw[fig,line width=0.5pt] ($(O)+(80.54:1.35)$) arc (80.54:22.11:1.35);
  \node[slab] at ($(O)+(109.01:1.80)$) {1};
  \node[slab] at ($(O)+(51.33:1.80)$)  {2};
\end{tikzpicture}
\figcap{3--27}
\end{bkfigure}}

% =============================================================== Fig 3-28
\newcommand{\FIGIIITWENTYEIGHT}{\begin{bkfigure}[0.55\linewidth]
\begin{tikzpicture}[scale=0.858]
  \coordinate (O) at (2.05,0.20);
  \draw[fig] (0.20,0.20) -- (4.15,0.20);
  \draw[fig] (O) -- (3.55,1.70);          % common side, at 45 degrees
  % both arcs run from a side of the angle to the common side, no gap
  \draw[fig,line width=0.5pt] ($(O)+(180:1.20)$) arc (180:45:1.20);
  \node[slab] at ($(O)+(112.5:0.85)$) {1};
  \draw[fig,line width=0.5pt] ($(O)+(0:0.67)$) arc (0:45:0.67);
  \node[slab] at ($(O)+(22.5:1.00)$) {2};
\end{tikzpicture}
\figcap{3--28}
\end{bkfigure}}

% ========================================================= Fig 3-29, 3-30
% 3-29: lines l and m crossing at O; the four regions cut by them, two of
% which are hatched (vertical / horizontal) so their overlap reads doubly hatched.
\newcommand{\FIGIIITWENTYNINETHIRTY}{\begin{bkfigure}[\linewidth]
\begin{minipage}{0.48\linewidth}\centering
\begin{tikzpicture}[scale=0.730]
  % The book's picture: A B C D at the corners of a rectangle, l = AC and
  % m = BD crossing at O.  VERTICAL rules shade the side of l that contains D;
  % HORIZONTAL rules shade the side of m that contains C.  The two families
  % therefore overlap in triangle OCD, and that doubly-shaded wedge IS the
  % interior of angle COD -- which is the whole point of the figure.
  \coordinate (B) at (0.35,2.55); \coordinate (Cc) at (3.95,2.55);
  \coordinate (A) at (0.35,0.35); \coordinate (D) at (3.95,0.35);
  \coordinate (O) at (intersection of A--Cc and B--D);
  % Each rule is drawn at its exact length rather than clipped: a clipped rule
  % still carries its full path in the PDF, so the collision audit would see
  % ink where the picture shows none.
  % D's side of l = triangle A D C: vertical rules from AD up to line AC
  \foreach \x in {0.45,0.55,...,3.90}
    \draw[fig,line width=0.35pt] (\x,0.35) -- (\x,{0.35+(\x-0.35)*0.611111});
  % C's side of m = triangle B C D: horizontal rules from line BD across to DC
  \foreach \y in {0.45,0.55,...,2.50}
    \draw[fig,line width=0.35pt] ({0.35+(2.55-\y)*1.636364},\y) -- (3.95,\y);
  % the two lines: carried past the far corner to take the label, and past the
  % near point as well -- in the book A and B are POINTS ON the lines, not ends
  \draw[fig] ($(A)!-0.07!(Cc)$) -- ($(A)!1.10!(Cc)$);
  \node[vlab,right] at ($(A)!1.10!(Cc)$) {$l$};
  \draw[fig] ($(B)!-0.07!(D)$) -- ($(B)!1.10!(D)$);
  \node[vlab,right] at ($(B)!1.10!(D)$) {$m$};
  \foreach \p in {A,B,Cc,D,O} {\dt{\p}}
  % A and B sit clear of the tails the two lines now carry past them
  \node[vlab] at ($(A)+(150:0.50)$) {$A$};  \node[vlab] at ($(B)+(95:0.58)$) {$B$};
  \node[vlab,above left] at (Cc) {$C$}; \node[vlab,below left] at (D) {$D$};
  \node[vlab] at ($(O)+(180:0.85)$) {$O$};   % the one unshaded wedge, as in the book
\end{tikzpicture}
\figcap{3--29}
\end{minipage}\hfill
\begin{minipage}{0.48\linewidth}\centering
\begin{tikzpicture}[scale=0.730]
  % r_1 lies on l_1 and r_2 on l_2, so each dashed piece is the ray OPPOSITE
  % its solid one: built as a reflection through O, never placed by eye.
  \coordinate (O) at (2.30,1.35);
  \coordinate (R1) at (4.30,2.55); \coordinate (R2) at (4.10,0.10);
  \coordinate (L1) at ($(O)!-0.90!(R1)$);   % opposite ray of r_1
  \coordinate (L2) at ($(O)!-1.00!(R2)$);   % opposite ray of r_2
  \draw[fig] (O) -- (R1); \draw[fig] (O) -- (R2);
  \draw[aux] (O) -- (L2); \draw[aux] (O) -- (L1);
  \node[vlab,right] at (R1) {$r_1$}; \node[vlab,right] at (R2) {$r_2$};
  \node[vlab,left] at (L2) {$l_2$};  \node[vlab,left] at (L1) {$l_1$};
  \dt{O} \node[vlab] at ($(O)+(90:0.62)$) {$O$};
\end{tikzpicture}
\figcap{3--30}
\end{minipage}
\end{bkfigure}}

% =============================================================== Fig 3-31
\newcommand{\FIGIIITHIRTYONE}{\begin{bkfigure}[0.70\linewidth]
\begin{tikzpicture}[scale=0.806]
  \coordinate (O) at (2.05,1.35);
  \coordinate (R1) at (4.05,2.45); \coordinate (R2) at (3.95,0.15);
  \coordinate (L1) at ($(O)!-0.85!(R1)$);   % opposite ray of r_1
  \coordinate (L2) at ($(O)!-0.95!(R2)$);   % opposite ray of r_2
  \draw[fig] (O) -- (R1); \node[vlab,right] at (R1) {$r_1$};
  \draw[fig] (O) -- (R2); \node[vlab,right] at (R2) {$r_2$};
  % The book draws l_1 and l_2 SOLID here, unlike Fig 3-30 where the rays that
  % are not sides of the angle are dashed -- checked on scans/scan04.pdf p20.
  \draw[fig] (O) -- (L2); \node[vlab,left] at (L2) {$l_2$};
  \draw[fig] (O) -- (L1); \node[vlab,left] at (L1) {$l_1$};
  \dt{O} \node[vlab] at ($(O)+(90:0.62)$) {$O$};
  % the separate segment PQ with A between; it runs a little past Q
  \coordinate (P) at (5.15,0.60); \coordinate (Q) at (7.15,0.60);
  \draw[fig] (P) -- ($(P)!1.14!(Q)$);
  \coordinate (Aa) at ($(P)!0.5!(Q)$);
  \foreach \p in {P,Q,Aa} {\dt{\p}}
  \node[vlab,above] at (P) {$P$}; \node[vlab,above] at (Q) {$Q$};
  \node[vlab,above] at (Aa) {$A$};
\end{tikzpicture}
\figcap{3--31}
\end{bkfigure}}

% =============================================================== Fig 3-32
\newcommand{\FIGIIITHIRTYTWO}{\begin{bkfigure}[0.60\linewidth]
\begin{tikzpicture}[scale=0.806]
  \node[vlab,left] at (0.20,2.45) {$r_2'$};
  \node[vlab,left] at (0.25,0.20) {$r_1'$};
  % the two full lines through O; r1' / r2' are the rays opposite r1 / r2
  \coordinate (ta) at (0.20,2.45); \coordinate (tb) at (4.05,0.30);
  \coordinate (ua) at (0.25,0.20); \coordinate (ub) at (4.15,2.35);
  \draw[fig] (ta) -- (tb);
  \draw[fig] (ua) -- (ub);
  \coordinate (O) at (intersection of ta--tb and ua--ub);
  \draw[fig] (O) -- (4.20,1.34); \node[vlab,right] at (4.20,1.34) {$r$};
  \node[vlab,right] at (4.15,2.35) {$r_1$}; \node[vlab,right] at (4.05,0.30) {$r_2$};
  \dt{O} \node[vlab] at ($(O)+(90:0.52)$) {$O$};
  % the two arcs mark angle(r_1,r) and angle(r,r_2): centred on O, ending on
  % the rays (O = (2.2524,1.3039); r_1 at 28.88, r at 1.06, r_2 at -29.18 deg)
  \draw[fig,line width=0.5pt] ($(O)+(28.88:1.40)$) arc (28.88:1.06:1.40);
  \draw[fig,line width=0.5pt] ($(O)+(1.06:1.40)$)  arc (1.06:-29.18:1.40);
\end{tikzpicture}
\figcap{3--32}
\end{bkfigure}}

% =============================================================== Fig 3-33
\newcommand{\FIGIIITHIRTYTHREE}{\begin{bkfigure}[0.72\linewidth]
\begin{tikzpicture}[scale=0.806]
  % l2 : the shallow line, carrying r2 to the right and r2' to the left
  \coordinate (L2R) at (5.75,1.95); \coordinate (L2L) at (0.30,1.75);
  \draw[fig] (L2L) -- (L2R);
  \node[vlab,right] at (L2R) {$l_2$}; \node[vlab,left] at (L2L) {$r_2'$};
  % l1 : the steep line; O is where they meet
  \coordinate (L1T) at (4.15,3.35); \coordinate (L1B) at (1.95,0.15);
  \draw[fig] (L1B) -- (L1T);
  \coordinate (O) at (intersection of L1B--L1T and L2L--L2R);
  \node[vlab,right] at (L1T) {$l_1$};
  \draw[aux] (L1B) -- (1.70,-0.20); \node[vlab,below] at (1.70,-0.20) {$r_1'$};
  \node[vlab,left]  at ($(O)!0.72!(L1T)$) {$r_1$};   % ray r1: above O
  \node[vlab,below] at ($(O)!0.80!(L2R)$) {$r_2$};   % ray r2: right of O
  % P lies in the interior of the angle; Q_3 is across BOTH lines from it, so
  % the book's segment Q_3 P has to meet l_1 in a point R and l_2 in a point S.
  % R and S are taken as intersections -- the proof's whole content.
  \coordinate (P) at (5.15,2.45);  \coordinate (Qc) at (1.85,1.15);
  \draw[fig] (Qc) -- (P);
  \coordinate (R) at (intersection of Qc--P and L1B--L1T);
  \coordinate (S) at (intersection of Qc--P and L2L--L2R);
  \foreach \p in {O,P,S,R} {\dt{\p}}
  \node[vlab] at ($(O)+(118.8:0.55)$) {$O$};
  \node[vlab,above right] at (P) {$P$};
  \node[vlab] at ($(S)+(60:0.75)$) {$S$};
  \node[vlab] at ($(R)+(308.5:0.48)$) {$R$};
  % the five sample positions.  Q_1, Q_3, Q_5 lie in the three regions the
  % lines cut off; Q_2 lies ON the ray r_2' and Q_4 ON the ray r_1' -- that
  % distinction is what Exercises *12-*14 turn on.
  \coordinate (Qa) at (2.55,3.05);                    % Q_1
  \coordinate (Qb) at ($(L2L)!0.22!(L2R)$);           % Q_2, on r_2'
  \coordinate (Qd) at ($(L1B)!0.12!(L1T)$);           % Q_4, on r_1'
  \coordinate (Qe) at (4.30,0.80);                    % Q_5
  \foreach \p in {Qa,Qb,Qc,Qd,Qe} {\dt{\p}}
  \node[vlab,above] at (Qa) {$Q_1$}; \node[vlab,above] at (Qb) {$Q_2$};
  \node[vlab] at ($(Qc)+(235:0.55)$) {$Q_3$};
  \node[vlab,right] at (Qd) {$Q_4$}; \node[vlab,above] at (Qe) {$Q_5$};
\end{tikzpicture}
\figcap{3--33}
\end{bkfigure}}

% ========================================================= Fig 3-34, 3-35
\newcommand{\FIGIIITHIRTYFOURTHIRTYFIVE}{\begin{bkfigure}[\linewidth]
\begin{minipage}{0.48\linewidth}\centering
\begin{tikzpicture}[scale=0.749]
  \coordinate (B) at (1.75,2.35); \coordinate (A) at (0.20,0.25);
  \coordinate (C) at (3.75,0.25);
  \draw[fig] (A)--(B)--(C)--cycle;
  \node[vlab,above] at (B) {$B$}; \node[vlab,left] at (A) {$A$};
  \node[vlab,right] at (C) {$C$};
\end{tikzpicture}
\figcap{3--34}
\end{minipage}\hfill
\begin{minipage}{0.48\linewidth}\centering
\begin{tikzpicture}[scale=0.749]
  \coordinate (B) at (1.55,2.35); \coordinate (A) at (0.20,0.25);
  \coordinate (C) at (3.55,0.25);
  \draw[fig] (A)--(B)--(C)--cycle;
  \coordinate (P) at (1.45,1.05); \coordinate (Q) at (3.55,2.20);
  \draw[fig] (P) -- (Q);
  \dt{P} \dt{Q}
  \node[vlab,above] at (B) {$B$}; \node[vlab,left] at (A) {$A$};
  \node[vlab,right] at (C) {$C$};
  \node[vlab,above,yshift=2pt] at (P) {$P$}; \node[vlab,above right] at (Q) {$Q$};
\end{tikzpicture}
\figcap{3--35}
\end{minipage}
\end{bkfigure}}

% =============================================================== Fig 3-36
\newcommand{\FIGIIITHIRTYSIX}{\begin{bkfigure}[0.50\linewidth]
\begin{tikzpicture}[scale=0.806]
  \coordinate (O) at (0.75,0.35); \coordinate (A) at (3.85,0.15);
  \coordinate (Ot) at (2.55,2.85); \coordinate (At) at (1.55,2.75);
  \draw[fig,dashed,dash pattern=on 3pt off 2.4pt] ($(O)!-0.23!(A)$) -- (O);
  \draw[fig] (O) -- (A);
  \draw[fig] (O) -- (Ot);
  \draw[fig] (A) -- (At);
  % B is where the ray from O meets the other side of the angle at A
  \coordinate (B) at (intersection of O--Ot and A--At);
  \dt{O} \dt{A} \dt{B}
  \node[vlab,above left] at (O) {$O$}; \node[vlab,right] at (A) {$A$};
  \node[vlab] at ($(B)+(93:0.80)$) {$B$};
\end{tikzpicture}
\figcap{3--36}
\end{bkfigure}}

% =============================================================== Fig 3-37
\newcommand{\FIGIIITHIRTYSEVEN}{\begin{bkfigure}[0.52\linewidth]
\begin{tikzpicture}[scale=0.806]
  \coordinate (V) at (1.35,0.65);
  \coordinate (M1) at (3.35,2.75);
  \draw[fig] (V) -- (3.95,0.65); \node[vlab,right] at (3.95,0.65) {$l_1$};
  \draw[fig] (V) -- (M1); \node[vlab,above] at (M1) {$m_1$};
  % l and m are the REST of the two lines: exact opposite rays through V
  \draw[fig,dashed,dash pattern=on 3pt off 2.4pt] (V) -- (0.15,0.65);
  \node[vlab,left] at (0.15,0.65) {$l$};
  \coordinate (Mo) at ($(V)!-0.45!(M1)$);
  \draw[fig,dashed,dash pattern=on 3pt off 2.4pt] (V) -- (Mo);
  \node[vlab,below] at (Mo) {$m$};
  \draw[fig,line width=0.5pt] ($(V)+(0:1.10)$) arc (0:46.40:1.10);
\end{tikzpicture}
\figcap{3--37}
\end{bkfigure}}

% ========================================================= Fig 3-38, 3-39
\newcommand{\FIGIIITHIRTYEIGHTTHIRTYNINE}{\begin{bkfigure}[\linewidth]
\begin{minipage}{0.48\linewidth}\centering
\begin{tikzpicture}[scale=0.675]
  % sides of the angle solid, the opposite rays r_1' and r_2' dashed
  \coordinate (l2L) at (0.30,1.05); \coordinate (l2R) at (4.55,1.25);
  \coordinate (ka) at (1.65,0.15);  \coordinate (kb) at (3.35,2.85);
  \coordinate (O) at (intersection of l2L--l2R and ka--kb);
  \draw[fig] (O) -- (l2R); \node[vlab,right] at (l2R) {$l_2$};
  \draw[aux] (O) -- (l2L); \node[vlab,left]  at (l2L) {$r_2'$};
  \draw[fig] (O) -- (kb);  \node[vlab,right] at (kb) {$l_1$};
  \draw[aux] (O) -- ($(O)!1.18!(ka)$);
  \node[vlab,below] at ($(O)!1.18!(ka)$) {$r_1'$};
  \node[vlab,below] at ($(O)!0.62!(l2R)$) {$r_2$};
  \node[vlab] at ($($(O)!0.78!(kb)$)+(147.8:0.42)$) {$r_1$};
  % Q_1 is on the same side of l_2 as r_1, the far side of l_1 from r_2;
  % R is where segment Q_1 P meets l_1 -- it has to land on the ray r_1.
  \coordinate (P) at (4.05,2.10); \coordinate (Q) at (1.35,2.45);
  \draw[fig] (Q) -- (P);
  \coordinate (R) at (intersection of Q--P and ka--kb);
  \foreach \p in {P,Q,R} {\dt{\p}}
  \node[vlab,right] at (P) {$P$}; \node[vlab,left] at (Q) {$Q_1$};
  \node[vlab] at ($(R)+(299.5:0.55)$) {$R$};
\end{tikzpicture}
\figcap{3--38}
\end{minipage}\hfill
\begin{minipage}{0.48\linewidth}\centering
\begin{tikzpicture}[scale=0.675]
  \coordinate (l2L) at (0.30,1.05); \coordinate (l2R) at (4.55,1.25);
  \coordinate (ka) at (1.65,0.15);  \coordinate (kb) at (3.35,2.85);
  \coordinate (O) at (intersection of l2L--l2R and ka--kb);
  \draw[fig] (O) -- (l2R); \node[vlab,right] at (l2R) {$l_2$};
  \draw[aux] (O) -- (l2L); \node[vlab,left]  at (l2L) {$r_2'$};
  \draw[fig] (O) -- (kb);  \node[vlab,right] at (kb) {$l_1$};
  \draw[aux] (O) -- ($(O)!1.18!(ka)$);
  \node[vlab,below] at ($(O)!1.18!(ka)$) {$r_1'$};
  \node[vlab,below] at ($(O)!0.62!(l2R)$) {$r_2$};
  \node[vlab] at ($($(O)!0.78!(kb)$)+(147.8:0.42)$) {$r_1$};
  % the one thing that separates this figure from 3-38: Q_2 lies ON the ray
  % r_2' itself, not merely on that side of it
  \coordinate (P) at (4.05,2.10); \coordinate (Q) at ($(l2L)!0.16!(l2R)$);
  \draw[fig] (Q) -- (P);
  \coordinate (R) at (intersection of Q--P and ka--kb);
  \foreach \p in {P,Q,R} {\dt{\p}}
  \node[vlab,right] at (P) {$P$}; \node[vlab,above] at (Q) {$Q_2$};
  \node[vlab] at ($(R)+(126:0.60)$) {$R$};
\end{tikzpicture}
\figcap{3--39}
\end{minipage}
\end{bkfigure}}


% ---- local macros (hoist into book preamble) ----
% Unnumbered star footnote, same reasoning as ch02's \IIstarnote: \starnote in
% brumfiel.sty emits a bare \footnotetext, which prints the footnote counter;
% the book marks all three of this chapter's notes with "*" and no number.
\newcommand{\IIIstarfn}[1]{\begingroup\renewcommand{\thefootnote}{}%
  \footnotetext{*#1}\endgroup}
% Starred theorems. The book prints "*Theorem 3-4." for the optional theorems
% of Section 3-6; this shares the `theorem' counter so numbering stays exact.
\theoremstyle{bkplain}
\newtheorem{IIIsthm}[theorem]{*Theorem}
% ---- end local macros ----

\begin{document}
\setcounter{chapter}{3}

\chapter*{\raggedright\Huge CONCERNING LINES}
\addcontentsline{toc}{chapter}{3.\ Concerning Lines}
\markboth{CONCERNING LINES}{}
\vspace{-1em}\noindent\rule{\linewidth}{0.6pt}\vspace{1em}
\figurekey

%========================================================= 3-1
\section*{3--1\quad INTRODUCTION}
\addcontentsline{toc}{section}{3--1\quad Introduction}

We have seen in Chapter 1 how from our everyday experience it is natural to
imagine the existence of things called points and lines. Even though one cannot
``see'' a point or a line, the study of such imaginary objects turns out to be
practically useful as well as interesting.

Because the objects of study in geometry are inventions of the mind, we must
state very carefully the properties we wish them to possess. Then if we start
with enough information about points and lines, other properties of points and
lines can be \emph{proved}. Many of these provable properties will be surprising,
and not at all obvious at the beginning. It is, of course, impossible to prove
anything about points and lines unless we agree in advance about some properties
that they are to have. These agreed-upon properties, or assumed properties, are
called \emph{postulates}, or \emph{axioms}.* We begin the proper study of plane
geometry in this chapter when we list our beginning postulates in Section 3--3.
The postulates should be simple enough to seem almost obvious and yet must be
sufficient for us to prove geometric facts that seem to be true in our everyday
life.%
\IIIstarfn{~The word ``postulate'' is derived from a Latin word meaning
``demand.'' ``Axiom'' comes from a Greek word meaning ``worthy.''}

A different method of procedure would be to list all the known geometric facts
and assume that they are true. Then when we had memorized all these facts, our
study of geometry would be completed. There are at least three objections to
such a plan. In the first place, we might not have listed all the facts needed
and we would have no experience in finding new ones. In the second place, it
would be a terrible strain upon the memory. To avoid memorizing everything, we
need to study how different facts are related and how the more complicated ones
depend upon other simpler ones. When we prove theorems from postulates or other
theorems, we are studying the way the geometric facts are related to each other;
this helps us to remember them. In the third place, we would lose the great
amount of pleasure that we can get out of proving from very simple and obvious
assumptions that certain other things are true. It is because of this pleasure
that geometry has been so attractive to men since the time of Euclid. The fact
that all geometry can be established from a small list of postulates gives to
geometry a kind of beauty that has appealed to men as different as Abraham
Lincoln and the Greek philosopher Plato.

%========================================================= 3-2
\section*{3--2\quad STRAIGHTNESS}
\addcontentsline{toc}{section}{3--2\quad Straightness}

Soon, when we list our postulates, we shall be talking about \emph{lines}, and of
course what we have in mind are things which are like the familiar \emph{straight}
lines of our experience. But we will call them simply \emph{``lines.''} There is
really nothing to be gained by putting in the extra word ``straight'' because it
would add nothing to our use of the postulates. It will turn out much later that
our lines \emph{are} straight in the sense that any path between two points,
other than a line, would be longer than the line.

When we draw pictures of lines we will, of course, draw them to \emph{appear}
straight, for this will help us see how the points and lines are related. But we
will not be concerned with the straightness of the lines because the figures are
meant to be only a guide to the proofs and to help us see what is going on.
\emph{The proofs are all completely independent of the figures.} By this we mean
that we must not say something is true because of the appearance of a figure;
rather, we must give reasons for its truth. The things that are true in our
geometry are only those things which can be proved by our use of our assumptions
and laws of logic. But as we progress in our study of geometry we will gradually
use figures more and more freely when it is quite clear how to fill in a complete
argument.

To understand why it is that we are not concerned with straightness at all, at
least to begin with, we shall look at a plane in two different ways. First,
imagine a sheet of paper with lines drawn on it so that some of the lines might
look like Fig.~3--1. Now suppose you have a pane of uneven glass between your
eyes and the paper, so that when you look at the drawing it looks like
Fig.~3--2. Note that in looking at the drawing, with or without the glass, you
are seeing the same lines. In one case they appear straight; in the other case
they do not. In the first figure $A$ and $B$ may appear to be the same distance
apart as $A'$ and $B'$, but this may not be so in the second figure. Therefore,
\emph{if it were possible to see the plane only through the glass, you would have
to be told which sets of points are ``straight'' lines, and whether or not $A$
and $B$ are the same distance apart as $A'$ and $B'$.}

\FIGIIIONETWO

In our study of geometry we shall be talking about imaginary things---points and
lines. We cannot really \emph{see} our lines. Hence we shall not call them
``straight'' but refer to them only as \emph{lines}. All the properties that we
need to know about these lines will be stated clearly in the postulates.

%========================================================= 3-3
\section*{3--3\quad THE POSTULATES}
\addcontentsline{toc}{section}{3--3\quad The postulates}

We imagine that there is a \emph{set} (or ``bunch,'' or ``collection'') of things
which we call \emph{points}. This set of points will be called \emph{the plane}.
In the plane certain other sets of points are called \emph{lines}. We represent
points by capital letters and lines by small letters. When a point $P$ is one of
the points of the line $l$ (Fig.~3--3), we say that ``$P$ is on $l$'' and also
that the line ``goes through $P$,'' or ``passes through $P$,'' or ``$l$ contains
$P$.'' If two lines $l$ and $m$ have a point $P$ in common (Fig.~3--4), that is,
if $P$ is on both $l$ and $m$, then we say that the lines \emph{intersect}.
Sometimes we will label two points with the same letter but different
\emph{subscripts}, such as $P_1$ and $P_2$. We read $P_1$ as ``$P$ sub-one'' and
$P_2$ as ``$P$ sub-two.''

\FIGIIITHREEFOURFIVE

The postulates will be divided into six groups, of which only the first group
will be given in this section. We shall give the other groups as we need them.
It is extremely important at the beginning of your study of geometry that you
understand the relationship between our mathematical postulates and the real
world. We decide \emph{what} we are going to assume about the imaginary points
and lines of our geometry by \emph{looking} at the real world. As a matter of
fact, in order to choose the postulates for our geometry, we need only to look at
some simple drawings. For example, the postulates in our first group are
suggested by Fig.~3--5. Study this figure to see if you can decide what these
postulates will be.

Of course, after we have decided what postulates to use, we must prove our
theorems by \emph{using these postulates and not by looking at other figures}.

\medskip
\noindent\textbf{Postulate Group I.}\quad This group contains three postulates,
indicated by I--1, I--2, and I--3.

\medskip
\noindent I--1 \emph{The three-point postulate}. There are at least three points
which do not all lie on one line.

\medskip
\noindent I--2 \emph{The point-line postulate}. For any two different points
there is exactly one line containing these points.

\medskip
\noindent I--3 \emph{The line-point postulate}. Every line contains at least two
points.

\medskip
The point-line postulate tells us one way we want lines to behave. The
line-point postulate \emph{starts} us on the way to getting points on a line. The
three-point postulate tells us that not all points are on one line. It gets us
off the line.

With only these postulates we cannot as yet prove very much. The following two
theorems are about all that can be formulated.

\begin{theorem}
Two different lines do not intersect in more than one point.
\end{theorem}

\noindent\emph{Proof.} The theorem is either true or false. We shall prove that
it cannot be false. Hence it must be true. We therefore suppose that the lines
intersect in at least two points and then show that this is impossible. Suppose
the lines are $l$ and $m$ and that they intersect in the two points $A$ and $B$
(Fig.~3--6). We would then have two different lines, $l$ and $m$, through $A$ and
$B$. This contradicts the point-line postulate. Therefore our assumption that
$l$ and $m$ intersect in more than one point is false and our theorem is true.

\begin{remark}
Several ideas from logic are involved in this proof. Comment on them.
\end{remark}

\begin{theorem}
There are at least three lines in the plane.
\end{theorem}

\FIGIIISIXSEVEN

\noindent\emph{Proof.}
\begin{steps}
\step{There are at least three points in the plane not all on one line. Call
      them $A$, $B$, and $C$.}{Three-point postulate.}
\step{Let $l$ be the line through $A$ and $B$, $m$ the line through $B$ and $C$,
      and $n$ the line through $C$ and $A$.}{Point-line postulate.}
\step{The three lines $l$, $m$, and $n$ are all different from one
      another.}{If any two were the same line, all three points $A$, $B$, and
      $C$ would be on one line, contrary to the way they were chosen.}
\end{steps}

The three points and lines are shown in Fig.~3--7. Of course, it looks just like
the drawing (Fig.~3--5) that suggested our first postulates.

\begin{remark}
At this point we cannot prove from our postulates that there are any more points
or lines than shown in Fig.~3--7.
\end{remark}

\exercisegroup{3--1}
% multicols dropped for this group: it carries figures, and the book sets each
% one immediately after the exercise that cites it, at the full measure.
\begin{exlist}
\item Restate the point-line and the three-point postulates in other words.
\item Prove that if $l$ is any line, there is at least one point not on $l$.
\item Explain why both $a$ and $b$ in Fig.~3--8 cannot represent lines if $P$ and
      $Q$ represent points.
\exfig{\FIGIIIEIGHT}
\item Prove that if $l$ is a line and $P$ is on $l$ then there is another point
      on $l$.
\item Why can we not prove there are more than three lines in our plane?
\item State what ``intersect'' means.
\item Suppose that $l$ is a line. List all the things you can say about this line
      and points of the plane. Give your reasons.
\item Given Fig.~3--9, what can you conclude about $l$ and $m$? Why?
\exfig{\FIGIIININE}
\item A line $l$ contains the points $P$ and $Q$. The line $m$ also contains $P$
      and $Q$. What can you conclude?
\item Two lines have a point $A$ in common. What can you say about the lines?
\item Prove that there are at least two lines on each point.
\item Why, when looking at Fig.~3--10, can we not prove that there are at least
      five points in the plane?
\exfig{\FIGIIITEN}
\end{exlist}

%========================================================= 3-4
\section*{3--4\quad POSTULATES OF BETWEENNESS}
\addcontentsline{toc}{section}{3--4\quad Postulates of betweenness}

We come now to our assumptions about the order, or arrangement, of points on a
line. It is interesting that Euclid did not state these assumptions, but used
them, without comment, wherever necessary. It was the German mathematician
Moritz Pasch (pronounced ``posh'') who first, in the latter part of the 19th
century, made a study of the properties of betweenness actually needed in
geometry. The postulates given below are a modification of those of Pasch. You
will have observed that the postulates of Group I seemed very obvious, and
perhaps those of Group II will also. But remember that we are trying to state
clearly what we imagine to be true about a plane, and hence we must be careful to
state \emph{everything} we are going to use.

\medskip
\noindent\textbf{Postulate Group II. Betweenness.}\quad There is a relation,
called \emph{between}, such that one of three points is sometimes between the
other two. Of course, if we consider \emph{any} three points, it does not
\emph{have} to happen that any one is between the others. But if for three points
$A$, $B$, and $C$ it is true that $B$ is between $A$ and $C$, we shall indicate
this by writing ``$ABC$'' or ``$CBA$,'' that is, with $B$ in the middle.

\medskip
\noindent II--1 \emph{The ABC betweenness postulate}. If $ABC$, that is, if $B$
is between $A$ and $C$, then $A$, $B$, and $C$ are three different (distinct)
points on a line. (A drawing of the line and points might look like Fig.~3--11.)

\FIGIIIELEVENTWELVE

\medskip
\noindent II--2 \emph{The three-point betweenness postulate}. For every three
points on a line, exactly one of them is between the other two. That is, if $R$,
$S$, and $T$ are on one line, then exactly one of the statements $RST$, $STR$,
and $TRS$ is true.

\medskip
\noindent II--3 \emph{The four-point betweenness postulate}. Any four points on a
line may be labeled $A_1$, $A_2$, $A_3$, $A_4$ so that the only betweenness
relations are the same as the order of the subscript numbers. That is, the only
betweenness relations are $A_1A_2A_3$, $A_1A_2A_4$, $A_1A_3A_4$, $A_2A_3A_4$. Of
course these may all be reversed as $A_3A_2A_1$, $A_4A_2A_1$, $A_4A_3A_1$,
$A_4A_3A_2$. A similar statement is to hold for any five points on a line, and in
fact for any finite number of points on a line.

\medskip
\noindent II--4 \emph{The line-building postulate}. If $A$ and $B$ are two
points, then there is at least one point $C$ such that $ABC$, and at least one
point $D$ such that $ADB$. (See Fig.~3--12.)

\medskip
Before stating the last of the betweenness postulates, we need some definitions.

\begin{definition}
Two points not on a line $l$ are said to be \emph{on the same side of $l$} if no
point of $l$ is between the two points. (See Fig.~3--13.) Two points not on a
line $l$ are said to be \emph{on opposite sides of $l$} if there is a point of
$l$ between the two points. (See Fig.~3--14.)
\end{definition}

\FIGIIITHIRTEENFOURTEEN

\medskip
\noindent II--5 \emph{The two-side postulate}. Each line separates the plane. By
this we mean that all points of the plane not on the line are divided into two
sets, called the two sides of the line, such that if $P$ and $Q$ belong to one of
these sets, then $P$ and $Q$ are on the same side of the line, whereas if $P$
belongs to one set and $Q$ belongs to the other, then $P$ and $Q$ are on opposite
sides of the line. (Briefly, this says that a line has two sides.)

\medskip
Before using these postulates, let us first discuss their meaning in everyday
life. If three men are standing in a line (Fig.~3--15) and $A$ cannot see $C$
because $B$ is in the way, then we say that $B$ is between $A$ and $C$. On the
other hand, $B$ and $C$ can see each other perfectly. Hence $A$ is not between
$B$ and $C$. That is, we do not have $BAC$.

\FIGIIIFIFTEEN

The $ABC$ betweenness postulate says that to talk about betweenness, we must have
the men ``lined up.'' The three-point betweenness postulate says that when three
men are lined up, exactly one of them is between the other two.

The four-point betweenness postulate gives a definite meaning to what we expect
to see when we say, ``Line up these four people.'' It tells exactly who is
between whom. An easy way to remember the content of this postulate is to observe
that of four points on a line, two of them are never between any others. These
are called the \emph{endpoints}. The other two points are between points exactly
twice.

The line-building postulate tells us first that there is always a point $C$
``beyond'' $B$ from $A$. In fact, Euclid stated (vaguely) something very much
like this as a postulate. But our postulate also tells us that there are points
between any two. Euclid never stated this. It is this postulate that gives the
existence of more than two points on a line, and therefore it is called the
line-building postulate. You might think about the relation between this
postulate and the real world. When you draw pictures of points, can you always
mark a picture of a third point between any two?

The two-side postulate says that a line has two sides. It merely expresses the
way we tell whether or not two of us are on the same or different sides of a
fence. If, when we look at each other, we see the fence between us, we know we
are on opposite sides of it. If we do not see the fence between us, we are on the
same side.

Now, using these postulates, we can define some familiar geometric objects.

\begin{definition}
If $A$ and $B$ are any two points, then the set of points consisting of $A$ and
$B$ and all the points between $A$ and $B$ will be called the \emph{line segment}
(or \emph{segment}) determined by $A$ and $B$. The segment will be indicated by
``$AB$.'' The points $A$ and $B$ are called the \emph{ends} of the segment. In
Fig.~3--16, $P$ is a point of the segment $AB$.
\end{definition}

\FIGIIISIXTEENSEVENTEEN

\begin{definition}
If $A$, $B$, and $C$ are three points not all on one line,* the set of points of
the segments $AB$, $BC$, $AC$ is called a \emph{triangle} (Fig.~3--17). We call
the points $A$, $B$, and $C$ \emph{vertices} of the triangle and the segments
$AB$, $AC$, $BC$ \emph{sides} of the triangle. To indicate ``triangle $ABC$,''
we write ``$\tri{ABC}$.''
\end{definition}
\IIIstarfn{~Points all lying on a line are said to be \emph{collinear}. Hence in
the definition we could have said, ``If $A$, $B$, and $C$ are not collinear
\ldots{}''.}

Before working the exercises, read all the postulates again to be sure that you
understand them.

\exercisegroup{3--2}
% multicols dropped for this group: it carries Fig. 3-18, and the book sets the
% figure immediately after the exercise that cites it, at the full measure.
\begin{exlist}
\item Draw a line and mark five points on it. Write out all the betweenness
      relations.
\item If $ABC$ and $ACD$, prove that the four points $A$, $B$, $C$, and $D$ lie
      on one line.
\item Give in your own words the definitions of ``segment,'' ``triangle,'' and
      ``two points are on the same side of a line.''
\item If $A$, $B$, $C$, and $D$ are four points of a line and we know that $CAD$
      and $ADB$, what are the other betweenness relations? Draw a picture. Is the
      picture proof?
\item[\stex 5.] If $A$, $B$, $C$, and $D$ are four points of a line and we know
      that $CAD$ and $CAB$, what are the other betweenness relations?
\item[6.] If $E$, $F$, $G$, and $H$ are four points on a line and we know that
      $EFG$, how many possible positions are there for $H$ that can be described
      in terms of our betweenness relations?
\end{exlist}

\begin{exlist}[start=7]
\item $A$ and $B$ are points on a line $l$. Write three true statements about
      points either on $l$ or not on $l$. Give reasons, as below.
\end{exlist}

\begin{steps}
\step{There is a point $C$ such that $ACB$.}{The line-building postulate.}
\step{\rule{0.92\linewidth}{0.4pt}}{\rule{0.92\linewidth}{0.4pt}}
\step{\rule{0.92\linewidth}{0.4pt}}{\rule{0.92\linewidth}{0.4pt}}
\step{\rule{0.92\linewidth}{0.4pt}}{\rule{0.92\linewidth}{0.4pt}}
\end{steps}

\begin{exlist}[start=8]
\item If $RST$ and $VRS$, relabel the points $R$, $S$, $T$, and $V$ as $A_1$,
      $A_2$, $A_3$, and $A_4$ (in some order) in accordance with II--3. If
      $CAD$, $BDE$, and $BCA$, relabel the points $A_1$, $A_2$, $A_3$, $A_4$, and
      $A_5$ in accordance with II--3.
\item Prove that there are at least four lines on each point.
\item If $DEF$ and $EGF$, then the four-point betweenness postulate implies that
      \rule{2.2em}{0.4pt} and \rule{2.6em}{0.4pt}.
\item If $RST$, then what postulate tells us that $R$, $S$, and $T$ are distinct
      points on a line?
\item[\stex 12.] Draw a line and choose two points on it, $A$ and $B$ (Postulate
      I--3). Mark a point $C$ such that $ABC$ (Postulate II--4). Mark a point $D$
      such that $BCD$ (Postulate II--4). Is it true that $ACD$? Why? Use II--3 to
      prove that $ACD$.
\item[\stex 13.] Draw a line and choose two points on it, $A$ and $B$. Choose a
      point $C$ such that $ACB$ (Postulate II--4). Now choose a point $D$ such
      that $CDB$. Is it true that $ACD$? that $ADB$? Use II--3 to prove this.
      Now choose a point $E$ such that $DEB$. What are all the betweenness
      relations between these five points? Can this process be indefinitely
      repeated? Describe, in your own words, the meaning of ``this process.'' How
      many points are there on a line segment?
\item[\stex 14.] Draw a line and choose a point $A$ on one side and a point $B$ on
      the other side. Now mark a point $C$ on the opposite side of the line from
      $A$. Points $B$ and $C$ will be on the same side of the line. How does this
      follow from II--5?
\item[\stex 15.] A line $l$ (Fig.~3--18) does not contain any of the vertices of
      a triangle. The line $l$ does intersect the side $AB$, but does not
      intersect the side $BC$. Prove that $l$ must intersect $AC$.
\exfig{\FIGIIIEIGHTEEN}
\item[\stex 16.] Draw a line and choose two points $A$ and $B$ on one side of the
      line. Choose $C$ so that $ACB$. Can you prove that $C$ and $A$ are on the
      same side of the line?
\item[\stex 17.] If $A$ is on line $l$ and $B$ is not on $l$ and $ACB$, prove
      that $C$ and $B$ are on the same side of $l$.
\item[\stex 18.] Prove that for every line $l$ there is a point on each side of
      $l$.
\end{exlist}

%========================================================= 3-5
\section*{3--5\quad RAYS AND ANGLES}
\addcontentsline{toc}{section}{3--5\quad Rays and angles}

It is a familiar fact from experience that a point $O$ on a line $l$ divides the
other points of the line into two sets---those on one side and those on the other
side of the point $O$ on the line $l$. (See Fig.~3--19.) From our postulates on
betweenness we can now make precise definitions of this observation. First we
need to know that a point on a line separates the line.

\FIGIIININETEEN

\begin{theorem}
{\normalfont The two-side theorem.} Let $O$ be any point of a line $l$. The
points of $l$, different from $O$, are separated into two sets such that:
(1) If $P$ and $Q$ are in the same set, then either $OPQ$ or $OQP$.
(2) If $P$ is in one set and $Q$ is in the other set, then $POQ$.
\end{theorem}

\noindent*\emph{Proof.} Choose a point $A$ not on $l$ and consider the line $m$
through $O$ and $A$ (see Fig.~3--20). We are using I--1 and I--2. By the two-side
postulate, II--5, the line $m$ separates the plane into the two sides of $m$. Let
one of the sets, which we seek, consist of the points of $l$ which are on one
side of $m$, and let the other set consist of the points of $l$ on the other side
of $m$.%
\IIIstarfn{~Starred proofs are difficult and may be omitted. However, the meaning
of the theorems must be understood.}

If $P$ and $Q$ are in one of these sets, they are by definition on the same side
of $m$, and therefore $O$ is not between them. Hence by II--2 we must have either
$OPQ$ or $OQP$. This proves (1).

Now suppose that $P$ is in one set and $Q$ in the other set (see Fig.~3--21).
Then by definition of the sets, $P$ and $Q$ are on different sides of $m$. By the
two-side postulate, II--5, there is a point of $m$ between $P$ and $Q$. But $l$
and $m$ have only point $O$ in common. Hence the point between must be $O$, and
we have $POQ$. This proves (2).

\FIGIIITWENTYTWENTYONE

\begin{remark}
The two sets of the theorem are not dependent on how we choose the line $m$,
because we end up with the same betweenness relations.
\end{remark}

\begin{definition}
The two sets of Theorem 3--3 are called the two \emph{sides} of $O$ on the line
$l$.
\end{definition}

\begin{definition}
A \emph{ray} from $O$, or a ray with end $O$, is a set of points consisting of
$O$ and all the points on one side of $O$ of a line $l$ through $O$ (see
Fig.~3--22).
\end{definition}

If $P$ is any point of the ray, we may speak of ``ray $OP$.'' Note that we may
describe the ray as the set of points consisting of $O$ and $P$ and all points
$Q$ such that $OQP$ and all points $R$ such that $OPR$. Refer back to
Theorem 3--3.

\FIGIIITWENTYTWOTWENTYTHREE

\exercise{3--3}
Consider a line $l$ and a ray $r$, not on $l$, from a point $P$ on $l$, as in
Fig.~3--23. What do we mean by ``$r$'s side of $l$''?

\begin{definition}
The set of all points on two rays from a point is called an \emph{angle}. The
point is called the \emph{vertex} of the angle, and the two rays are called the
\emph{sides} of the angle. If the two rays are on the same line, the angle is
called a \emph{straight angle}. If $O$ is the vertex and $A$ and $B$ are points
other than $O$ on the sides of the angle, we shall speak of ``angle $AOB$'' and
write this as ``$\ang{AOB}$.'' (See Fig.~3--24.)
\end{definition}

\FIGIIITWENTYFOURTWENTYFIVE

\begin{remark}
Note that what we call the angle is the set of points in the \emph{sides} of the
angle. Points ``between'' the rays are \emph{not} part of the angle. Later we
discuss the inside and the outside of an angle. It will be convenient to have a
way to clearly indicate the sides of the angles. This can be done as in
Fig.~3--25(a), and the angle can be numbered. Since the arrowheads indicating the
sides are not vital, we can as well use the simpler notation shown in
Fig.~3--25(b) and (c). \emph{Remember that this notation is just a way of
indicating the sides of the angle.}
\end{remark}

\begin{definition}
If two lines $l$ and $m$ intersect at $O$ and the four rays from $O$ are as
indicated in Fig.~3--26, then angles 1 and 2 are said to be \emph{vertical} with
respect to each other. For brevity, we say that 1 and 2 are a pair of vertical
angles. Angles 3 and 4 are also vertical to each other.
\end{definition}

\FIGIIITWENTYSIX

\begin{definition}
If two angles have a common vertex and a common side and if the other sides are
on opposite sides of the line on the common side, then the angles are said to be
\emph{adjacent} to each other (Fig.~3--27).
\end{definition}

\FIGIIITWENTYSEVEN

\begin{definition}
Two adjacent angles whose noncommon sides form a straight angle are said to be
\emph{supplementary and adjacent} to each other (Fig.~3--28).
\end{definition}

\FIGIIITWENTYEIGHT

\exercisegroup{3--4}
\begin{multicols}{2}
\begin{exlist}
\item State in your own words the definitions of ``ray,'' of ``angle,'' and of
      ``straight angle.''
\item How ``long'' are the sides of an angle?
\item If $\ang{AOB}$ and $\ang{DOC}$ are vertical to each other with $B$, $O$,
      $D$ collinear, what kind of angle is $\ang{AOC}$?
\item If $\ang{AOC}$ is a straight angle and $OR$ is a ray from $O$, how are
      angles $AOR$ and $ROC$ related?
\item Two lines intersect at $O$. With the four rays from $O$, how many different
      angles with vertex $O$ can be found? Pick out the pairs of vertical angles.
\item Three lines intersect at $O$. With the six rays from $O$, how many
      different angles with vertex $O$ can be formed? How many pairs of vertical
      angles are there?
\item Draw two angles with a common vertex but which are not adjacent to each
      other.
\item Draw two angles with a common vertex and side but which are not adjacent.
\item State Definitions 3--7, 3--8, and 3--9 in your own words.
\item If $\ang{AOB}$ is not a straight angle, how many angles are there that are
      vertical with respect to $\ang{AOB}$?
\item Why is it not sensible to say, ``$\ang{AOB}$ is a vertical angle''?
\end{exlist}
\end{multicols}

%========================================================= 3-6
\section*{3--6\quad INTERIOR AND EXTERIOR}
\addcontentsline{toc}{section}{3--6\quad Interior and exterior}

When we draw a picture of an angle (not a straight angle) on a sheet of paper, we
can tell by looking at it which points of the plane are inside and which are
outside the \emph{picture} of the angle. But in our development of geometry
starting from but a few postulates, we must \emph{define} what these sets are.
(Why?) Definition 3--10 is one that can be understood by a blind person provided
he knows what an angle is and that a line separates the plane. To appreciate the
definition, consider Fig.~3--29. The point $D$ lies on one side of $l$, which
side has been shaded with vertical lines. The point $C$ lies on one side of $m$,
which side has been shaded with horizontal lines. The points in the plane which
are shaded twice (i.e., where the vertical and horizontal lines cross) fill up
the \emph{interior} of $\ang{COD}$.

\FIGIIITWENTYNINETHIRTY

\exercisegroup{3--5}
\begin{multicols}{2}
\begin{exlist}
\item Copy the lines $l$ and $m$ and the points $A$, $B$, $C$, $D$ on a sheet of
      paper and shade all the points of the plane which are (a) on the same side
      of $l$ as $B$ and the same side of $m$ as $A$, (b) on the same side of $l$
      as $B$ and on the same side of $m$ as $C$, (c) on the same side of $l$ as
      $D$ and on the same side of $m$ as $A$.
\item In each case name the angle whose interior has been described.
\end{exlist}
\end{multicols}

\begin{definition}
Consider an angle, not a straight angle, with vertex $O$ and whose sides are the
rays $r_1$ and $r_2$ on the lines $l_1$ and $l_2$, respectively. The
\emph{interior} of the angle is the set of all points on the same side of $l_1$
as $r_2$ and on the same side of $l_2$ as $r_1$. The \emph{exterior} of the angle
is the set of all points in the plane which are not points of the angle and are
not in the interior of the angle. (Fig.~3--30.)
\end{definition}

There are several properties of the interior and exterior of angles that we can
now prove. These properties are all quite obvious when we look at a figure. But
even when something seems obvious from a picture, we must be sure that we have
the right postulates so that our geometry agrees with our feelings about what
should be true---that is, with the picture. It is really not at all obvious that
our \emph{geometry} agrees with our \emph{picture}. Let us see that this is so by
proving some ``obvious'' theorems.

In all these theorems we shall be concerned with an angle whose vertex is $O$,
which is not a straight angle, and whose sides are the rays $r_1$ and $r_2$.

\begin{IIIsthm}
If $P$ and $Q$ are both in the interior of $\ang{O}$, then the segment $PQ$ is in
the interior of $\ang{O}$.
\end{IIIsthm}

\noindent\emph{Proof.} We must show that if $A$ is any point between $P$ and $Q$,
then $A$ is in the interior of the angle (Fig.~3--31). This is easy if we use the
result of Exercise 16 in Exercise Group 3--2. By that exercise, since $P$ and $Q$
are on one side of $l_1$, then $A$ is also on that same side of $l_1$. Similarly,
$A$ is on the same side of $l_2$, as are both $P$ and $Q$. Hence, by the
definition of \emph{interior} of an angle, $A$ is in the interior of the angle
formed by rays $r_1$ and $r_2$.

\FIGIIITHIRTYONE

\begin{IIIsthm}
If $r$ is a ray from $O$ other than $r_1$ and $r_2$, then either all the points
of $r$ other than $O$ lie in the interior of $\ang{O}$ or they lie in the
exterior of $\ang{O}$.
\end{IIIsthm}

\noindent\emph{Proof.} Let $r_1'$ and $r_2'$ be the rays from $O$ on $l_1$ and
$l_2$ other than $r_1$ and $r_2$. (See Fig.~3--32.) If $r = r_1'$ or
$r = r_2'$, then $r$ is in the exterior. If $r$ is neither $r_1'$ nor $r_2'$,
then $r$ lies on a definite side of $l_1$ and a definite side of $l_2$ because of
the way a ray is defined. Hence $r$ is in either the interior or exterior of
$\ang{O}$.

\FIGIIITHIRTYTWO

\begin{IIIsthm}
If $P$ is in the interior of $\ang{O}$ and $Q$ is in the exterior of $\ang{O}$,
then the segment $PQ$ contains a point of the angle.
\end{IIIsthm}

\noindent\emph{Proof.} There are five possible locations of the point $Q$,
represented by the points $Q_1, \ldots, Q_5$ in Fig.~3--33. (Describe these
positions in terms of sides of lines.) For positions $Q_1$, $Q_2$, $Q_4$, $Q_5$
the proofs of the theorem are quite easy and are assigned as exercises below.
Hence we restrict our attention here to the position $Q_3$.

\FIGIIITHIRTYTHREE

Since $Q_3$ is on opposite sides of both $l_1$ and $l_2$ from $P$ (why?), the
segment $Q_3P$ must contain a point $R$ of $l_1$ and a point $S$ of $l_2$. We
assert that either $R$ or $S$ is a point of the angle, that is, a point of either
$r_1$ or $r_2$. If $R$ is on $r_1$, our segment contains a point of the angle.
Suppose that $R$ is not on $r_1$. Then $R$ is on $r_1'$ and hence has the same
relative position as $Q_4$. (Why?) But then the segment $RP$ contains a point of
the angle which must be $S$. (Why?) Since $Q_3RP$ and $RSP$, we have $Q_3SP$.
(Why?)

\begin{definition}
The \emph{interior} (or \emph{inside}) of a triangle is the set of all points
which are interior to each of the angles of the triangle. The \emph{exterior} of
a triangle is the set of points which are neither on the triangle nor in the
interior.
\end{definition}

If the triangle is $ABC$ (see Fig.~3--34), an alternative definition of interior
would be ``the set of all points on the same side of $AB$ as $C$, on the same
side of $BC$ as $A$, and on the same side of $CA$ as $B$.'' Show that these two
definitions say the same thing.

We can prove the following remark by using Theorem 3--6. The remark is an
``obvious'' theorem that is hard to prove. We will not need this result later.

\begin{remark}
If $P$ is in the interior of $\tri{ABC}$ and $Q$ is in the exterior, then the
segment $PQ$ contains a point of the triangle (Fig.~3--35).
\end{remark}

\FIGIIITHIRTYFOURTHIRTYFIVE

\exercisegroup{3--6}
% multicols dropped for this group: it carries figures, and the book sets each
% one immediately after the exercise that cites it, at the full measure.
\begin{exlist}
\item Referring to Fig.~3--36, draw an angle with vertex $O$. Pick a point $A$ on
      one side of $\ang{O}$. Draw an angle with vertex $A$ with one side the ray
      $AO$ and the other side a ray intersecting the other side of $\ang{O}$ in
      $B$. Shade the interior of $\ang{O}$ and the interior of $\ang{A}$. Are
      points in both interiors also in the interior of $\tri{OBA}$?
\exfig{\FIGIIITHIRTYSIX}
\item Draw a $\tri{ABC}$. If a point $P$ is on the same side of $AB$ as $C$ and
      on the same side of $BC$ as $A$, show that $P$ is not necessarily on the
      same side of $AC$ as $B$.
\item Draw two angles whose interiors have no points in common.
\item A straight angle does not have an interior. Tell why we do not try to
      define interior for a straight angle.
\item State in your own words the definitions of ``interior of an angle'' and
      ``interior of a triangle.''
\item For a $\tri{ABC}$, the point $P$ is in the interior of $\ang{A}$ and not in
      the interior of $\ang{B}$. Can you be sure that $P$ is in the exterior of
      $\tri{ABC}$?
\item Consider the angle of Fig.~3--37, with sides $l_1$ and $m_1$ indicated on
      lines $l$ and $m$. Tell where the point $P$ is if:
      \begin{itemize}[label=,leftmargin=1.1em,itemsep=0pt,topsep=1pt]
      \item (a) $P$ is on the opposite side of $m$ from $l_1$.
      \item (b) $P$ is on the opposite side of $l$ from $m_1$.
      \item (c) $P$ is on $l_1$.
      \item (d) $P$ is on $m_1$.
      \item (e) $P$ is on the same side of $l$ as $m_1$.
      \item (f) $P$ is on the same side of $l$ as $m_1$ and on the same side of
            $m$ as $l_1$.
      \end{itemize}
\exfig{\FIGIIITHIRTYSEVEN}
\item Prove that the following definition is equivalent to Definition 3--11: The
      interior of $\tri{ABC}$ is the set of all points in the interior of
      $\ang{A}$ and also in the interior of $\ang{B}$. (With this definition it
      is not too difficult to prove the truth of the remark preceding this group
      of exercises.)
\item In $\ang{AOB}$, not a straight angle, point $P$ is between $A$ and $B$.
      Prove that $P$ is in the interior of the angle.
\item In triangle $ABC$ we have the relations $AXB$, $AYC$, and $XZY$. Prove
      that $Z$ is in the interior of the triangle.
\item If $\ang{AOB}$ is not a straight angle and we have the relations $XOB$ and
      $XYA$, prove that $Y$ is in the exterior of $\ang{AOB}$.
\item[\stex 12.] In Theorem 3--6 suppose $Q$ is in position $Q_1$, that is, on
      the same side of $l_2$ as $r_1$ and on the opposite side of $l_1$ as $r_2$
      (see Fig.~3--38). Show that there is a point $R$ of $r_1$ between $Q$ and
      $P$. [\emph{Hint:} Certainly there is a point $R$ of $l_1$ between $Q$ and
      $P$. Show that it is impossible that $R$ be on $r_1'$.]
\exfig{\FIGIIITHIRTYEIGHT}
\item[\stex 13.] In Theorem 3--6 suppose $Q$ is in position $Q_2$, that is, on
      $r_2'$ (see Fig.~3--39). Show that there is a point $R$ of $r_1$ between
      $Q$ and $P$.
\exfig{\FIGIIITHIRTYNINE}
\item[\stex 14.] Tell why the positions $Q_4$ and $Q_5$ of Fig.~3--33 are much
      the same as positions $Q_2$ and $Q_1$, respectively. If you have done
      positions $Q_1$ and $Q_2$, why are positions $Q_5$ and $Q_4$ easy?
\end{exlist}

%========================================================= end matter
\section*{REVIEW OF CHAPTER 3}
\addcontentsline{toc}{section}{Review of Chapter 3}

\begin{exlist}
\item The fundamental relations of betweenness are expressed in the postulates.
      Study them again to be sure you know them.
\item Definitions of the following concepts were given in this chapter. Be sure
      that you know all of them.
\end{exlist}

\begin{multicols}{2}
\begin{itemize}[label=,leftmargin=1.2em,itemsep=0pt,topsep=2pt]
\item two sides of a line (see II--5)
\item segment
\item triangle
\item collinear (see footnote to Def.~3--3)
\item two sides of a point on a line
\item ray
\item angle
\item straight angle
\item vertical angles
\item adjacent angles
\item supplementary angles
\item interior of an angle
\item exterior of an angle
\item interior of a triangle
\item exterior of a triangle
\end{itemize}
\end{multicols}

\begin{exlist}[start=3]
\item The following statements have been proved in this chapter. First re-read
      them, and then study their proofs again.
      \begin{itemize}[label=,leftmargin=1.4em,itemsep=1pt,topsep=2pt]
      \item (a) Two different lines intersect in at most one point.
      \item (b) There are at least three lines.
      \item (c) A point $O$ on a line separates the line into two sides.
      \item (d) If $P$ and $Q$ are in the interior of $\ang{O}$, then the segment
            $PQ$ is in the interior of $\ang{O}$.
      \item (e) A ray from the vertex of an angle lies, except for the vertex,
            either in the interior or in the exterior of the angle.
      \item (f) If $P$ is in the interior of $\ang{O}$ and $Q$ is in the
            exterior, then the segment $PQ$ intersects the angle.
      \end{itemize}
\end{exlist}

\begin{center}\bfseries Important Remark\end{center}

By assuming a few properties of points and lines, we have been able to
\emph{define} many familiar types of figures and to \emph{prove} that these
figures have properties that may seem obvious. \emph{Therefore we have proved
that some of our figures agree with our geometry.} From now on in the book, we
will often determine properties of betweenness by just looking at figures. The
definitions and theorems of this chapter will enable us to use our figures in the
confident belief that we could always go back and supply proofs. Therefore, we
have completed the first stage in our development of geometry. The theorems
proved in this chapter need not be memorized because the facts are all familiar
from drawings. However, the postulates and definitions of this chapter are very
important, and therefore should be memorized.

\begin{center}\bfseries Review Questions\end{center}

\begin{exlist}
\item If $A$ and $B$ are distinct points, what can you say about a point $P$
      which is on the segment $AB$? Why?
\item What postulate did we use to define the sides of a point on a line?
\item A line $l$ contains the distinct points $R$ and $S$. A line $m$ also
      contains these points. What do you conclude?
\item If $ABC$, and $B$ is on line $l$ but $A$ is not on $l$, is the point $C$ on
      $l$? Why?
\item Tell whether the following statements are true or false, and give
      explanations for your answers.
      \begin{itemize}[label=,leftmargin=1.4em,itemsep=1pt,topsep=2pt]
      \item (a) For $\tri{ABC}$, if the point $S$ is in the interior of $\ang{A}$,
            $\ang{B}$, and $\ang{C}$, then it is a theorem that $S$ is in the
            interior of the triangle.
      \item (b) If $APB$, then $P$ belongs to the segment $AB$, by definition.
      \item (c) If $ARQ$ and $AGR$, then $AGQ$.
      \item (d) A point $P$ not on an angle must be in either the interior or
            exterior.
      \item (e) We postulate that a point on a line separates the line.
      \item (f) If $AQB$, then $A$ and $B$ are on a ray from $Q$.
      \item (g) A line segment contains a definite number of points.
      \end{itemize}
\item What does it mean to say that three or more points are noncollinear? How
      can we be sure that there are such sets of points?
\item If $ADB$ and $BAC$, what are the other betweenness relations?
\item Lines $s$ and $t$ intersect at $Q$ forming rays $s_1$, $s_2$ and $t_1$,
      $t_2$. Draw a figure and indicate the pairs of vertical angles. Which pairs
      of angles are adjacent? supplementary?
\item If $A$ and $B$ are on opposite sides of a line $l$, and $B$ and $C$ are on
      opposite sides of $l$, then $A$ and $C$ are on the same side of $l$. Why is
      this true?
\item Explain in your own words what it means for a point to be in the interior
      of an angle.
\end{exlist}

\begin{center}\bfseries Review Test\end{center}

\begin{exlist}
\item Define the following: (a) angle, (b) segment, (c) when two points are on
      the same side of a line, (d) triangle.
\item Tell whether true or false:
      \begin{itemize}[label=,leftmargin=1.4em,itemsep=1pt,topsep=2pt]
      \item (a) $A$, $B$, $C$ are three points on a line and $ABC$ and $CAB$.
      \item (b) If $PRU$, then $P$, $R$, and $U$ are distinct points on a line.
      \item (c) It is a theorem that a line separates the plane.
      \item (d) It is a theorem that a point separates a line.
      \item (e) It is a theorem that there is a point between any two points.
      \item (f) Betweenness is an undefined relation.
      \item (g) If $ABC$, then $CBA$.
      \item (h) If $E$, $F$, $G$ are not on $l$ and such that the segment $EF$
            does not intersect $l$, while the segment $FG$ does intersect $l$,
            then the segment $EG$ intersects $l$.
      \end{itemize}
\item State in your own words:
      \begin{itemize}[label=,leftmargin=1.4em,itemsep=1pt,topsep=2pt]
      \item (a) The point-line postulate.
      \item (b) The $ABC$ betweenness postulate.
      \item (c) The line-building postulate.
      \end{itemize}
\end{exlist}

\begin{center}\bfseries ALGEBRA REVIEW\end{center}

In this review you are supposed to use the properties of addition and
multiplication of numbers (both positive and negative), along with the
definitions of subtraction and division, to make simple proofs. All letters here
refer to rational numbers (fractions). These exercises show how definitions are
used in algebra.

\medskip
\noindent\textbf{Definition.}\quad If $a = b + x$, then $x$ is called ``$a$
\emph{minus} $b$'' and is written ``$x = a - b$.'' It is also called ``$b$
\emph{subtracted from} $a$.''

\begin{exlist}
\item If $c = p - q$, what is $p$? Why?
\item Why does $7 - 3 = 4$? Write $7 - 3 = x$, and ask what property $x$ must
      have.
\item Using the definition, what can you say about these numbers:
      $\frac{9}{5} - \frac{3}{5}$? $\frac{1}{2} - \frac{1}{3}$?
      $\frac{7}{8} - \frac{3}{4}$? $4.678 - 2.589$?
\end{exlist}

\medskip
\noindent\textbf{Definition.}\quad $-b$ (called ``\emph{negative} $b$'') is
defined to be $0 - b$. That is, $-b = 0 - b$ and is a number such that
$b + (-b) = 0$.

\begin{exlist}[start=4]
\item Why does $-7 = 0 - 7$? $-3 = 0 - 3$?
      $-\frac{1}{2} = 0 - \frac{1}{2}$?
\item Show that $9 + (-3) = 9 - 3$, $11 + (-8) = 11 - 8$,
      $\frac{5}{6} + (-\frac{1}{3}) = \frac{5}{6} - \frac{1}{3}$.
\item[\stex 6.] Prove that $a + (-b) = a - b$. [\emph{Hint:} Let $a - b = x$ so
      that $a = b + x$. Then $a + (-b) = b + x + (-b)$.]
\item[7.] Show that $-(7 + 3) = (-7) + (-3)$, $-(2 + 7) = (-2) + (-7)$.
\item[8.] Show that $-(a + b) = (-a) + (-b)$.
\end{exlist}

\medskip
\noindent\textbf{Definition.}\quad If $a = bx$ and $b \neq 0$, then $x$ is called
``$a$ \emph{divided by} $b$'' and is written ``$x = a \div b$'' or
``$x = a/b$.''

\begin{exlist}[start=9]
\item Explain why in this last definition we required that $b \neq 0$.
      [\emph{Hint:} Would $x$ be definite otherwise?]
\item According to the definition, $\frac{1}{2}$ is a number such that
      $1 = 2(\frac{1}{2})$. What then are the numbers $\frac{3}{4}$,
      $-\frac{2}{3}$, $4/{-5}$, $\frac{7}{16}$, $\frac{5}{10}$, $\frac{1}{1}$,
      $\frac{2}{2}$?
\item Why does $b \cdot a/b = a$?
\item Prove that $a/b + c/b = (a + c)/b$.
\item Show that $\frac{1}{2} = \frac{2}{4}$. [\emph{Hint:} Let
      $\frac{1}{2} = x$. Then $1 = 2 \cdot x$. Why? Hence $2 = 4x$. Why? Hence
      $\frac{2}{4} = x$. Why?]
\item Prove that if $b \neq 0$ and $c \neq 0$, then $a/b = ac/bc$.
\end{exlist}

\end{document}
